
\documentclass[11pt,a4paper,oneside,openany,fontset=fandol]{ctexbook}

\usepackage[a4paper,top=2.5cm,bottom=2.5cm,left=2.7cm,right=2.7cm,headheight=15pt]{geometry}
\usepackage{amsmath,amssymb,mathtools,bm}
\usepackage{booktabs,array,tabularx}
\usepackage{enumitem}
\usepackage{xparse}
\usepackage{tikz}
\usetikzlibrary{arrows.meta,positioning,calc,decorations.pathmorphing}
\usepackage[hidelinks]{hyperref}
\usepackage{microtype}
\usepackage{setspace}
\usepackage{caption}

\setlength{\parindent}{2em}
\setlength{\parskip}{0.15em}
\setlist{nosep,leftmargin=2.2em}
\linespread{1.12}
\pagestyle{plain}

\NewDocumentEnvironment{routebox}{}{
  \begin{quote}\small\noindent\textbf{主线}\par\medskip
}{
  \end{quote}
}
\NewDocumentEnvironment{keybox}{O{核心结论}}{
  \begin{quote}\small\noindent\textbf{#1}\par\medskip
}{
  \end{quote}
}
\NewDocumentEnvironment{remarkbox}{O{说明}}{
  \begin{quote}\small\noindent\textbf{#1}\par\medskip
}{
  \end{quote}
}
\NewDocumentEnvironment{derivationbox}{O{推导}}{
  \begin{quote}\small\noindent\textbf{#1}\par\medskip
}{
  \end{quote}
}

\definecolor{Ink}{gray}{0}
\definecolor{Navy}{gray}{0}
\definecolor{Teal}{gray}{0}
\definecolor{Gold}{gray}{0}
\definecolor{SoftBlue}{gray}{1}
\definecolor{SoftGold}{gray}{1}
\definecolor{SoftGray}{gray}{1}
\definecolor{DeepRed}{gray}{0}

\newcommand{\dd}{\mathrm{d}}
\newcommand{\ee}{\mathrm{e}}
\newcommand{\ii}{\mathrm{i}}
\newcommand{\Tr}{\operatorname{Tr}}
\newcommand{\Var}{\operatorname{Var}}
\newcommand{\diag}{\operatorname{diag}}
\newcommand{\ket}[1]{\left|#1\right\rangle}
\newcommand{\bra}[1]{\left\langle#1\right|}
\newcommand{\braket}[2]{\left\langle#1\,\middle|\,#2\right\rangle}
\newcommand{\avg}[1]{\left\langle #1\right\rangle}
\newcommand{\pb}[2]{\left\{#1,#2\right\}_{\mathrm{PB}}}
\newcommand{\vct}[1]{\bm{#1}}
\newcommand{\mat}[1]{\mathbf{#1}}
\newcommand{\order}{\mathcal{O}}
\newcommand{\Wig}{\mathrm{W}}

\begin{document}


\begin{titlepage}
\centering
\vspace*{4.5cm}
{\Huge\bfseries 化学中的数学\par}
\vspace{1cm}
{\large 刘剑老师课程笔记 \par}
\vspace{0.5cm}
{\large nanomsj \par}
\vspace{0.5cm}
{\large transformed to \LaTeX by gpt 5.6 Sol\par}
\vspace{1.5cm}
{\normalsize \today\par}
\vfill

\end{titlepage}

\frontmatter
\tableofcontents
\vfill
说明：
本笔记根据本人在刘剑老师 2026 年暑期课程中的手写课堂笔记整理而成。为便于阅读，手写内容经 GPT 辅助转换为 LaTeX。
受限于原始笔记的完整性，以及自动识别与转换过程可能产生的偏差（正文未经校对），文中或有错漏及排版问题。如发现任何问题，欢迎在 GitHub 上提交 Issue 或 pr ，亦可发送邮件至 msj@stu.pku.edu.cn 联系我修正。也欢迎通过邮件与我交流课程内容！
本笔记仅供学习与交流使用。
\mainmatter


\chapter{全局脉络}

\begin{routebox}
量子--经典联系有两条彼此呼应的路线：
\[
\hbar\to 0
\quad\Longleftrightarrow\quad
\text{相位快速振荡，驻相路径占主导}
\]
以及
\[
t=-\ii\hbar\beta
\quad\Longleftrightarrow\quad
\text{实时间传播子变为 Boltzmann 算符}.
\]
前者把路径积分与经典作用量联系起来，后者把量子统计力学化为虚时间路径积分，并进一步得到环形高分子同构。
\end{routebox}

\begin{center}
\begin{tikzpicture}[
  node distance=8mm and 11mm,
  box/.style={draw=Navy!75,rounded corners=2mm,fill=SoftBlue,
              minimum height=9mm,align=center,inner xsep=4mm},
  arr/.style={-{Latex[length=2.4mm]},thick,color=Gold}
]
\node[box] (ham) {Hamilton\\相空间};
\node[box,right=of ham] (lio) {Liouville\\体积守恒};
\node[box,right=of lio] (ens) {系综平均\\遍历性};
\node[box,below=of ham] (lag) {Lagrange\\作用量};
\node[box,right=of lag] (prop) {量子传播子\\虚时间};
\node[box,right=of prop] (tro) {Trotter\\短时拆分};
\node[box,below=of prop] (pi) {路径积分};
\node[box,left=of pi] (rp) {环形高分子\\PIMD};
\node[box,right=of pi] (wig) {Wigner\\相空间};

\draw[arr] (ham)--(lio);
\draw[arr] (lio)--(ens);
\draw[arr] (ham)--(lag);
\draw[arr] (lag)--(prop);
\draw[arr] (prop)--(tro);
\draw[arr] (tro)--(pi);
\draw[arr] (pi)--(rp);
\draw[arr] (pi)--(wig);
\end{tikzpicture}
\end{center}

\section{量子表象的并列视角}

量子力学可以采用多种等价表象：
\[
\text{矩阵表象},\quad
\text{波函数表象},\quad
\text{路径积分表象},\quad
\text{量子相空间表象},\quad
\text{量子轨迹表象}.
\]
本笔记以波函数表象为起点，经过传播子和路径积分，最后进入 Wigner 相空间表象。

\chapter{量子--经典联系与 Born--Oppenheimer 近似}

\section{$\hbar\to 0$ 的经典极限}

定态量子力学满足
\[
\hat H\ket{\phi_n}=E_n\ket{\phi_n}.
\]
在位置表象与动量表象中，
\[
\braket{\vct r}{\phi}=\phi(\vct r),\qquad
\braket{\vct p}{\phi}=\widetilde\phi(\vct p).
\]

经典极限并不是简单地把所有公式中的 $\hbar$ 机械地置零，而是指相对于体系的典型作用量尺度，$\hbar$ 足够小。在 Wigner--Weyl 对应下，
\[
\frac{1}{\ii\hbar}[\hat A,\hat B]
\;\longleftrightarrow\;
\pb{A_{\Wig}}{B_{\Wig}}+\order(\hbar^2),
\]
量子对易子在最低阶化为 Poisson 括号。路径积分中则表现为
\[
K\sim\int\mathcal D x\,\ee^{\ii S[x]/\hbar},
\]
当 $\hbar$ 很小时，不满足 $\delta S=0$ 的相邻路径因相位快速振荡而互相抵消，经典驻相路径占主导。

\begin{keybox}
量子到经典的核心桥梁可概括为
\[
\text{对易子}\to\text{Poisson 括号},
\qquad
\text{全体路径}\to\text{经典驻相路径}.
\]
\end{keybox}

\section{相同粒子的交换对称性}

对于两个相同粒子，交换粒子标签相当于交换坐标参数。玻色子波函数对称，
\[
\Psi(\vct r_1,\vct r_2)=\Psi(\vct r_2,\vct r_1),
\]
费米子波函数反对称，
\[
\Psi(\vct r_1,\vct r_2)=-\Psi(\vct r_2,\vct r_1).
\]
在二维空间中还允许任意子统计：
\[
\Psi(\vct r_1,\vct r_2)
=\ee^{\ii\theta}\Psi(\vct r_2,\vct r_1),
\qquad \theta\in[0,2\pi).
\]
其中 $\theta=0$ 对应玻色子，$\theta=\pi$ 对应费米子。某些模型可先在玻色型参数区域求解，再对统计参数作解析延拓；但这种做法依赖具体模型，并非无条件适用于所有体系。

\section{Born--Oppenheimer 近似与势能面}

电子质量远小于原子核质量，典型电子运动时间尺度快于核运动。Born--Oppenheimer（BO）近似先固定原子核坐标 $\{\vct R_I\}$，求解电子本征问题：
\[
\hat H_e\!\left(\{\vct r_i\};\{\vct R_I\}\right)
\Psi_n\!\left(\{\vct r_i\};\{\vct R_I\}\right)
=
E_n\!\left(\{\vct R_I\}\right)
\Psi_n\!\left(\{\vct r_i\};\{\vct R_I\}\right).
\]
在显式单位下，电子 Hamiltonian 可写为
\[
\hat H_e=
-\sum_i\frac{\hbar^2}{2m_e}\nabla_i^2
-\sum_{iI}\frac{Z_Ie^2}{4\pi\varepsilon_0|\vct r_i-\vct R_I|}
+\sum_{i<j}\frac{e^2}{4\pi\varepsilon_0|\vct r_i-\vct r_j|}.
\]
第 $n$ 个绝热势能面为
\[
U_n(\{\vct R_I\})
=
E_n(\{\vct R_I\})
+
\sum_{I<J}
\frac{Z_IZ_Je^2}{4\pi\varepsilon_0|\vct R_I-\vct R_J|}.
\]
若只保留基态电子面，核波函数满足
\[
\left[
-\sum_{I=1}^{M}\frac{\hbar^2}{2M_I}\nabla_I^2
+U_0(\{\vct R_I\})
\right]\Theta(\{\vct R_I\})
=E\Theta(\{\vct R_I\}).
\]

\begin{remarkbox}[势能面的含义]
$U_n(\{\vct R_I\})$ 是定义在 $3M$ 维核构型空间上的标量函数。经典分子动力学通常让原子核在某一势能面上运动；量子核运动则在同一势能面上求解核 Schrödinger 方程或进行路径积分采样。
\end{remarkbox}

\section{BO 近似的失效与核量子效应}

将总波函数展开为
\[
\Psi(\vct r,\vct R)
=
\sum_n\Theta_n(\vct R)\Psi_n(\vct r;\vct R)
\]
后，核动能会产生非绝热耦合
\[
\vct d_{mn}^{(I)}(\vct R)
=
\left\langle
\Psi_m(\vct r;\vct R)
\middle|
\nabla_I
\Psi_n(\vct r;\vct R)
\right\rangle .
\]
当不同电子态能量接近、势能面发生避免交叉或锥形交叉时，这些耦合不可忽略，电子与原子核不能再分离。光异构化中的无辐射跃迁是典型例子。

即使 BO 分离成立，原子核也未必可以完全经典化。零点能、隧穿和量子离域会影响氢键、振动谱与同位素效应；常温水的 Raman 光谱便需要考虑显著的核量子效应。

\chapter{Hamilton 力学与相空间}

\section{Hamilton 方程}

设广义坐标与共轭动量分别为 $\vct x$、$\vct p$。一般的 Hamiltonian 可写成
\[
H(\vct x,\vct p)
=
\frac12\vct p^{\mathsf T}\mat M^{-1}(\vct x)\vct p
+V(\vct x).
\]
Hamilton 方程为
\[
\dot{\vct x}=\frac{\partial H}{\partial\vct p},
\qquad
\dot{\vct p}=-\frac{\partial H}{\partial\vct x}.
\]
因此
\[
\dot{\vct x}=\mat M^{-1}(\vct x)\vct p,
\]
而第 $a$ 个动量分量满足
\[
\dot p_a
=
-\frac12\vct p^{\mathsf T}
\frac{\partial\mat M^{-1}}{\partial x_a}
\vct p
-\frac{\partial V}{\partial x_a}.
\]
若质量矩阵为常量，
\[
\dot{\vct x}=\mat M^{-1}\vct p,
\qquad
\dot{\vct p}=-\nabla V(\vct x).
\]

\section{能量守恒}

沿 Hamilton 轨道，
\begin{align}
\frac{\dd H}{\dd t}
&=
\frac{\partial H}{\partial\vct x}\cdot\dot{\vct x}
+
\frac{\partial H}{\partial\vct p}\cdot\dot{\vct p}
+
\frac{\partial H}{\partial t}\\
&=
\frac{\partial H}{\partial\vct x}\cdot
\frac{\partial H}{\partial\vct p}
-
\frac{\partial H}{\partial\vct p}\cdot
\frac{\partial H}{\partial\vct x}
+
\frac{\partial H}{\partial t}
=
\frac{\partial H}{\partial t}.
\end{align}
因此，若 $H$ 不显含时间，则 $\dd H/\dd t=0$。

\section{相空间与 Liouville 定理}

位置和动量共同组成相空间：
\[
\vct z=
\begin{pmatrix}
\vct x\\
\vct p
\end{pmatrix},
\qquad
\dot{\vct z}
=
\mat J\nabla_{\vct z}H,
\qquad
\mat J=
\begin{pmatrix}
0&I\\
-I&0
\end{pmatrix}.
\]
由初值 $\vct z_0$ 经过时间 $t$ 演化到 $\vct z_t$，定义切映射
\[
\mat M_t=\frac{\partial\vct z_t}{\partial\vct z_0}.
\]
相空间体积元的变化为
\[
\dd\vct z_t
=
\left|\det\mat M_t\right|\dd\vct z_0.
\]

\begin{derivationbox}[Liouville 定理的行列式证明]
利用 Jacobi 公式，
\[
\frac{\dd}{\dd t}\ln\det\mat M_t
=
\Tr\!\left(
\mat M_t^{-1}\dot{\mat M}_t
\right)
=
\nabla_{\vct z}\cdot\dot{\vct z}.
\]
Hamilton 流的散度为
\begin{align}
\nabla_{\vct z}\cdot\dot{\vct z}
&=
\sum_a
\left(
\frac{\partial\dot x_a}{\partial x_a}
+
\frac{\partial\dot p_a}{\partial p_a}
\right)\\
&=
\sum_a
\left(
\frac{\partial^2H}{\partial x_a\partial p_a}
-
\frac{\partial^2H}{\partial p_a\partial x_a}
\right)
=0.
\end{align}
所以 $\det\mat M_t$ 为常数。初始时 $\mat M_0=I$，故
\[
\boxed{\det\mat M_t=1}.
\]
即 Hamilton 动力学保持相空间体积。
\end{derivationbox}

等价地，Hamilton 流是辛变换：
\[
\mat M_t^{\mathsf T}\mat J\mat M_t=\mat J.
\]

\chapter{系综、遍历性与统计平均}

\section{微正则与正则系综}

固定能量 $E$ 的微正则分布为
\[
\rho_E(\vct z)
=
\frac{\delta(E-H(\vct z))}{\Omega(E)},
\qquad
\Omega(E)
=
\int\dd\vct z\,\delta(E-H(\vct z)).
\]
正则分布为
\[
\rho_\beta(\vct z)
=
\frac{\ee^{-\beta H(\vct z)}}{Z_{\mathrm{cl}}},
\qquad
Z_{\mathrm{cl}}
=
\int\dd\vct z\,\ee^{-\beta H(\vct z)}
\]
（量纲化时还需除以适当的 $h^F$，见第 \ref{sec:classicalZ} 节）。可观测量 $B$ 的正则平均为
\[
\avg{B}_\beta
=
\frac{\int\dd\vct z\,
\ee^{-\beta H(\vct z)}B(\vct z)}
{\int\dd\vct z\,
\ee^{-\beta H(\vct z)}}.
\]

\section{遍历性假设}

对一条充分长的轨道 $\vct z_t$，时间平均定义为
\[
\overline B_T
=
\frac1T\int_0^T\dd t\,B(\vct z_t).
\]
若在给定能量面上满足遍历性，则
\[
\lim_{T\to\infty}\overline B_T
=
\int\dd\vct z\,\rho_E(\vct z)B(\vct z)
\equiv\avg{B}_E.
\]

\begin{remarkbox}[一个容易混淆的点]
不带恒温器的 Hamilton 动力学保持能量，因而在理想遍历条件下采样的是微正则分布，而不是正则分布。正则平均可通过恒温分子动力学、Langevin 动力学或 Monte Carlo 获得。
\end{remarkbox}

正则积分可以分解为能量壳上的微正则积分：
\begin{align}
\int\dd\vct z\,\ee^{-\beta H(\vct z)}B(\vct z)
&=
\int\dd E\,\ee^{-\beta E}
\int\dd\vct z\,
\delta(E-H(\vct z))B(\vct z)\\
&=
\int\dd E\,\ee^{-\beta E}\Omega(E)\avg{B}_E.
\end{align}
这正是“在每个能量壳上遍历，再对能量作 Boltzmann 加权”的严格形式。

\section{涨落与统计误差}

可观测量的方差为
\[
\Var(B)
=
\avg{B^2}-\avg{B}^2,
\]
标准差为
\[
\Delta B=\sqrt{\Var(B)}.
\]
数值采样中，相邻时间点通常相关，故有效独立样本数小于总步数；误差分析需要考虑自相关时间。

\section{高维积分与采样}

基本 Gaussian 积分
\[
\int_{-\infty}^{\infty}\ee^{-a x^2+b x}\dd x
=
\sqrt{\frac{\pi}{a}}\,
\ee^{b^2/(4a)},
\qquad \operatorname{Re}a>0,
\]
在动量积分、传播子和虚时间路径积分中反复出现。

高维空间若逐点网格化，计算量随自由度指数增长。Monte Carlo 直接按概率密度抽样，分子动力学则生成连续轨道；二者都以“重要性采样”替代对整个高维空间的穷举。

\chapter{Lagrange 力学与作用量}

\section{作用量与 Euler--Lagrange 方程}

给定端点
\[
x(0)=x_0,\qquad x(t_f)=x_f,
\]
作用量定义为
\[
S[x]
=
\int_0^{t_f}\dd t\,L(x,\dot x,t).
\]
对路径作变分 $x(t)\to x(t)+\delta x(t)$，并令端点变分为零：
\[
\delta x(0)=\delta x(t_f)=0.
\]
则
\begin{align}
\delta S
&=
\int_0^{t_f}\dd t
\left(
\frac{\partial L}{\partial x}\delta x
+
\frac{\partial L}{\partial\dot x}\delta\dot x
\right)\\
&=
\int_0^{t_f}\dd t
\left(
\frac{\partial L}{\partial x}
-
\frac{\dd}{\dd t}
\frac{\partial L}{\partial\dot x}
\right)\delta x
+
\left.
\frac{\partial L}{\partial\dot x}\delta x
\right|_0^{t_f}.
\end{align}
边界项为零。由于内部的 $\delta x(t)$ 任意，作用量驻值条件 $\delta S=0$ 给出
\[
\boxed{
\frac{\dd}{\dd t}
\frac{\partial L}{\partial\dot x}
-
\frac{\partial L}{\partial x}
=0
}.
\]

对于
\[
L=T-V=\frac12m\dot x^2-V(x),
\]
有
\[
m\ddot x=-\frac{\partial V}{\partial x},
\]
即 Newton 方程。

\section{Lagrange 与 Hamilton 描述的关系}

定义共轭动量
\[
p=\frac{\partial L}{\partial\dot x},
\]
并作 Legendre 变换
\[
H(x,p)=p\dot x-L(x,\dot x).
\]
由此得到 Hamilton 方程。两种描述完全等价，但侧重点不同：

\begin{center}
\begin{tabular}{lll}
\toprule
描述 & 基本变量 & 核心结构\\
\midrule
Lagrange & $(x,\dot x)$ & 路径、作用量、边值问题\\
Hamilton & $(x,p)$ & 相空间、辛结构、初值问题\\
\bottomrule
\end{tabular}
\end{center}

\chapter{量子传播子与虚时间}

\section{时间演化算符}

含时 Schrödinger 方程为
\[
\ii\hbar\frac{\dd}{\dd t}\ket{\psi(t)}
=
\hat H\ket{\psi(t)}.
\]
若 $\hat H$ 不显含时间，
\[
\ket{\psi(t)}
=
\ee^{-\ii\hat Ht/\hbar}\ket{\psi(0)}.
\]
时间演化算符
\[
\hat U(t)=\ee^{-\ii\hat Ht/\hbar}
\]
作用在任意初态上。

位置本征态满足
\[
\hat x\ket{x}=x\ket{x},
\qquad
\braket{x}{x'}=\delta(x-x'),
\qquad
\int\dd x\,\ket{x}\bra{x}=1.
\]
因此
\[
\psi(y,t)
=
\int\dd x\,
\underbrace{\bra{y}\ee^{-\ii\hat Ht/\hbar}\ket{x}}_{K(y,x;t)}
\psi(x,0),
\]
其中
\[
K(y,x;t)
=
\bra{y}\ee^{-\ii\hat Ht/\hbar}\ket{x}
\]
称为传播子。

\section{动量本征态与 Fourier 归一化}

位置表象中
\[
\hat p=-\ii\hbar\frac{\partial}{\partial x}.
\]
动量本征方程
\[
\hat p\ket{p}=p\ket{p}
\]
给出
\[
-\ii\hbar\frac{\partial}{\partial x}\braket{x}{p}
=
p\braket{x}{p},
\]
故
\[
\braket{x}{p}=C\ee^{\ii px/\hbar}.
\]
利用
\[
\delta(x-y)
=
\int\dd p\,\braket{x}{p}\braket{p}{y}
\]
和 Fourier 表示
\[
\delta(x-y)
=
\frac{1}{2\pi\hbar}
\int_{-\infty}^{\infty}
\dd p\,
\ee^{\ii p(x-y)/\hbar},
\]
得到
\[
\boxed{
\braket{x}{p}
=
\frac{1}{\sqrt{2\pi\hbar}}
\ee^{\ii px/\hbar}
}.
\]

\section{自由粒子传播子}

自由粒子 Hamiltonian 为
\[
\hat H_0=\frac{\hat p^2}{2m}.
\]
在传播子中插入动量完备关系：
\begin{align}
K_0(y,x;t)
&=
\int\dd p\,
\braket{y}{p}
\ee^{-\ii p^2t/(2m\hbar)}
\braket{p}{x}\\
&=
\frac{1}{2\pi\hbar}
\int\dd p\,
\exp\left[
\frac{\ii}{\hbar}p(y-x)
-\frac{\ii t}{2m\hbar}p^2
\right].
\end{align}
完成平方并使用 Gaussian 积分，得到
\[
\boxed{
K_0(y,x;t)
=
\sqrt{\frac{m}{2\pi\ii\hbar t}}
\exp\left[
\frac{\ii m(y-x)^2}{2\hbar t}
\right]
}.
\]

自由粒子的经典路径为
\[
x_{\mathrm{cl}}(t')
=
x+\frac{y-x}{t}t',
\]
其经典作用量
\[
S_{\mathrm{cl}}(y,x;t)
=
\int_0^t\dd t'\,\frac12m\dot x_{\mathrm{cl}}^2
=
\frac{m(y-x)^2}{2t}.
\]
因此
\[
K_0(y,x;t)
=
\sqrt{\frac{m}{2\pi\ii\hbar t}}
\ee^{\ii S_{\mathrm{cl}}/\hbar}.
\]

\section{虚时间与 Boltzmann 算符}

令
\[
t=-\ii\hbar\beta,
\]
则
\[
\ee^{-\ii\hat Ht/\hbar}
=
\ee^{-\beta\hat H}.
\]
自由粒子的虚时间核为
\[
\boxed{
\bra{y}\ee^{-\beta\hat H_0}\ket{x}
=
\sqrt{\frac{m}{2\pi\hbar^2\beta}}
\exp\left[
-\frac{m(y-x)^2}{2\hbar^2\beta}
\right]
}.
\]
实时间中的振荡相位变成正定的 Gaussian 权重，这使量子统计力学可以转写为概率采样问题。

\begin{remarkbox}[谐振子的精确虚时间核]
对于
\[
\hat H=\frac{\hat p^2}{2m}+\frac12m\omega^2\hat x^2,
\]
有
\[
\bra{x}\ee^{-\beta\hat H}\ket{x'}
=
\sqrt{\frac{m\omega}{2\pi\hbar\sinh(\beta\hbar\omega)}}
\exp\left[
-\frac{m\omega}{2\hbar\sinh(\beta\hbar\omega)}
\left(
(x^2+x'^2)\cosh(\beta\hbar\omega)-2xx'
\right)
\right].
\]
\end{remarkbox}

\chapter{Trotter 拆分与路径积分}

\section{插入完备关系}

将总时间分成 $P$ 段，
\[
t=P\Delta t.
\]
传播子可写为
\[
K(x_P,x_0;t)
=
\int\prod_{j=1}^{P-1}\dd x_j
\prod_{j=0}^{P-1}
\bra{x_{j+1}}
\ee^{-\ii\hat H\Delta t/\hbar}
\ket{x_j}.
\]
这一步把一条连续路径离散为 $P$ 个短时片段，内部插入了 $P-1$ 个位置完备关系。

\section{Trotter--Suzuki 拆分}

设
\[
\hat H=\hat T+\hat V,
\qquad
[\hat T,\hat V]\ne 0.
\]
一阶拆分为
\[
\ee^{-\epsilon(\hat T+\hat V)}
=
\ee^{-\epsilon\hat T}
\ee^{-\epsilon\hat V}
+\order(\epsilon^2),
\]
对称二阶拆分为
\[
\boxed{
\ee^{-\epsilon(\hat T+\hat V)}
=
\ee^{-\epsilon\hat V/2}
\ee^{-\epsilon\hat T}
\ee^{-\epsilon\hat V/2}
+\order(\epsilon^3)
}.
\]
单步误差为 $\order(\epsilon^3)$，固定总时间后全局误差为 $\order(P^{-2})$。

\section{短时传播核}

对一维体系
\[
\hat H=\frac{\hat p^2}{2m}+V(\hat x),
\]
令 $\Delta t=t/P$，对称拆分给出
\begin{align}
\bra{x_{j+1}}
\ee^{-\ii\hat H\Delta t/\hbar}
\ket{x_j}
&\approx
\exp\left[
-\frac{\ii\Delta t}{2\hbar}
\left(V(x_{j+1})+V(x_j)\right)
\right]\\
&\quad\times
\sqrt{\frac{m}{2\pi\ii\hbar\Delta t}}
\exp\left[
\frac{\ii m(x_{j+1}-x_j)^2}{2\hbar\Delta t}
\right].
\end{align}

代入全传播子并令 $P\to\infty$，
\[
K(x_f,x_i;t)
=
\int_{x(0)=x_i}^{x(t)=x_f}
\mathcal D x\,
\exp\left[
\frac{\ii}{\hbar}
\int_0^t\dd t'\,
\left(
\frac12m\dot x^2-V(x)
\right)
\right].
\]
这就是 Feynman 路径积分：
\[
\boxed{
K=\int\mathcal D x\,\ee^{\ii S[x]/\hbar}
}.
\]

\section{虚时间路径积分}

在虚时间中令
\[
\beta_P=\frac{\beta}{P},
\]
短时密度矩阵为
\begin{align}
\bra{x_{j+1}}
\ee^{-\beta_P\hat H}
\ket{x_j}
&\approx
\left(
\frac{m}{2\pi\hbar^2\beta_P}
\right)^{1/2}\\
&\quad\times
\exp\left[
-\frac{m(x_{j+1}-x_j)^2}{2\hbar^2\beta_P}
-\frac{\beta_P}{2}
\left(V(x_{j+1})+V(x_j)\right)
\right].
\end{align}
连续极限可写为
\[
\rho(x_f,x_i;\beta)
=
\int_{x(0)=x_i}^{x(\beta\hbar)=x_f}
\mathcal D x(\tau)\,
\ee^{-S_E[x]/\hbar},
\]
其中 Euclidean 作用量为
\[
S_E[x]
=
\int_0^{\beta\hbar}\dd\tau
\left[
\frac12m\left(\frac{\dd x}{\dd\tau}\right)^2+V(x)
\right].
\]
求迹时端点相接：
\[
x(\beta\hbar)=x(0),
\]
因此热路径是虚时间上的闭合环。

\chapter{环形高分子同构与 PIMD}

\section{量子配分函数的离散路径积分}

量子配分函数
\[
Z=\Tr\ee^{-\beta\hat H}
=
\int\dd x_1\,
\bra{x_1}\ee^{-\beta\hat H}\ket{x_1}.
\]
插入 $P-1$ 个位置完备关系，并采用对称 Trotter 拆分：
\begin{align}
Z_P
&=
\left(
\frac{mP}{2\pi\beta\hbar^2}
\right)^{P/2}
\int\prod_{j=1}^{P}\dd x_j\\
&\quad\times
\exp\left[
-\frac{\beta}{P}\sum_{j=1}^{P}V(x_j)
-\frac{mP}{2\beta\hbar^2}
\sum_{j=1}^{P}(x_j-x_{j+1})^2
\right],
\end{align}
其中
\[
x_{P+1}=x_1,
\qquad
Z=\lim_{P\to\infty}Z_P.
\]

定义
\[
\beta_P=\frac{\beta}{P},
\qquad
\omega_P=\frac{1}{\beta_P\hbar}
=\frac{P}{\beta\hbar},
\]
则
\[
Z_P
=
C_P
\int\dd\vct x\,
\exp\left[
-\beta_P
\sum_{j=1}^{P}
\left(
\frac12m\omega_P^2(x_j-x_{j+1})^2+V(x_j)
\right)
\right].
\]

\begin{keybox}[环形高分子同构]
一个量子粒子的平衡配分函数，等价于一个由 $P$ 个经典“珠子”构成的环形高分子：
\[
\boxed{
U_P(\vct x)
=
\sum_{j=1}^{P}
\left[
\frac12m\omega_P^2(x_j-x_{j+1})^2+V(x_j)
\right]
}.
\]
相邻珠子由谐弹簧连接，每个珠子都受到原始势能 $V$。
\end{keybox}

\begin{center}
\begin{tikzpicture}[scale=1.1]
\def\n{10}
\def\r{1.65}
\foreach \i in {1,...,\n}{
  \coordinate (P\i) at ({360/\n*(\i-1)}:\r);
}
\foreach \i in {1,...,\n}{
  \pgfmathtruncatemacro{\j}{mod(\i,\n)+1}
  \draw[decorate,decoration={coil,aspect=0.45,segment length=3.5pt,amplitude=2.2pt},
        color=Gold,thick] (P\i)--(P\j);
}
\foreach \i in {1,...,\n}{
  \filldraw[fill=SoftBlue,draw=Navy,thick] (P\i) circle (2.6pt);
}
\node[align=center,color=Ink] at (0,0) {虚时间闭合\\$x_{P+1}=x_1$};
\node[color=Teal] at (0,-2.35) {每个珠子感受 $V(x_j)$，相邻珠子由谐弹簧耦合};
\end{tikzpicture}
\end{center}

\section{多维形式}

对 $F$ 个自由度，质量矩阵为 $\mat M$，则
\begin{align}
Z_P
&=
C_P
\int\prod_{j=1}^{P}\dd\vct x_j\,
\exp\left[-\beta_P U_P(\vct x_1,\ldots,\vct x_P)\right],\\
U_P
&=
\sum_{j=1}^{P}
\left[
\frac12\omega_P^2
(\vct x_j-\vct x_{j+1})^{\mathsf T}
\mat M
(\vct x_j-\vct x_{j+1})
+
V(\vct x_j)
\right].
\end{align}
归一化常数为
\[
C_P=
\left[
\det(\mat M)
\right]^{P/2}
\left(
\frac{P}{2\pi\beta\hbar^2}
\right)^{FP/2}.
\]

\section{引入虚拟动量}

Gaussian 恒等式允许在构型积分中引入任意正定的虚拟质量矩阵 $\widetilde{\mat M}$：
\[
\int\dd\vct x\,\ee^{-\beta V_{\mathrm{eff}}(\vct x)}
=
\left(\frac{\beta}{2\pi}\right)^{D/2}
\frac{1}{\sqrt{\det\widetilde{\mat M}}}
\int\dd\vct x\,\dd\vct p\,
\ee^{-\beta H_{\mathrm{eff}}(\vct x,\vct p)},
\]
其中
\[
H_{\mathrm{eff}}
=
\frac12\vct p^{\mathsf T}
\widetilde{\mat M}^{-1}\vct p
+
V_{\mathrm{eff}}(\vct x).
\]
对环形高分子，常用
\[
H_P
=
\sum_{j=1}^{P}
\frac12\vct p_j^{\mathsf T}
\widetilde{\mat M}_j^{-1}\vct p_j
+
U_P(\vct x_1,\ldots,\vct x_P),
\]
并在温度 $\beta_P^{-1}$ 下采样
\[
\ee^{-\beta_P H_P}.
\]

\begin{remarkbox}[虚拟动量不是物理动量]
$\vct p_j$ 与 $\widetilde{\mat M}_j$ 是为构造高效采样动力学而引入的辅助变量。改变虚拟质量会改变采样效率和人工动力学频率，但不改变平衡构型分布。
\end{remarkbox}

这正是路径积分分子动力学（PIMD）的基础：用经典 Hamilton 方程推进一个扩展的环形高分子体系，从而采样量子 Boltzmann 分布。

\section{位置型可观测量的估计量}

对于只依赖位置的算符
\[
\hat B=B(\hat x),
\]
量子热平均为
\[
\avg{\hat B}
=
\frac{\Tr(\ee^{-\beta\hat H}\hat B)}
{\Tr(\ee^{-\beta\hat H})}.
\]
由于迹具有循环不变性，每个珠子地位等价，得到原始估计量
\[
\boxed{
B_P(\vct x)
=
\frac1P\sum_{j=1}^{P}B(x_j)
}.
\]
特别地，环形高分子的质心
\[
\bar x=\frac1P\sum_{j=1}^{P}x_j
\]
是重要的集体坐标，但一般有
\[
\frac1P\sum_jB(x_j)\ne B(\bar x)
\]
（除非 $B$ 为线性函数）。

\section{动量算符的导数表示}

对任意算符 $\hat A$，
\[
\bra{x}\hat A\hat p\ket{x'}
=
\ii\hbar\frac{\partial}{\partial x'}
\bra{x}\hat A\ket{x'}.
\]
令 $\hat A=\ee^{-\beta_P\hat H}$，并采用原始短时核
\[
K_P(x,x')
\approx
\left(
\frac{m}{2\pi\hbar^2\beta_P}
\right)^{1/2}
\exp\left[
-\frac{m(x-x')^2}{2\hbar^2\beta_P}
-\frac{\beta_P}{2}\left(V(x)+V(x')\right)
\right],
\]
则
\[
\frac{\bra{x}\ee^{-\beta_P\hat H}\hat p\ket{x'}}
{K_P(x,x')}
\approx
\ii\left[
\frac{m(x-x')}{\beta_P\hbar}
-\frac{\beta_P\hbar}{2}V'(x')
\right].
\]
对无磁场、实势能的平衡体系，时间反演对称性给出
\[
\avg{\hat p}=0.
\]
对于动能或 $\hat p^2$，实际计算通常采用热力学估计量或 virial 估计量，以降低方差。

\section{经典配分函数与相空间量子化}
\label{sec:classicalZ}

经典配分函数需要用 Planck 常数给相空间体积定标：
\[
Z_{\mathrm{cl}}
=
\frac{1}{h^F}
\int\dd\vct x\,\dd\vct p\,
\ee^{-\beta H(\vct x,\vct p)},
\qquad h=2\pi\hbar.
\]
对于一维无限深势阱（长度 $L$），
\[
E_n=\frac{\hbar^2\pi^2n^2}{2mL^2}.
\]
在 $L\to\infty$ 的连续极限中，量子求和变为动量积分：
\[
Z
=
\sum_{n=1}^{\infty}\ee^{-\beta E_n}
\longrightarrow
\frac{L}{h}
\int_{-\infty}^{\infty}
\dd p\,\ee^{-\beta p^2/(2m)}
=
\frac{L}{h}
\sqrt{\frac{2\pi m}{\beta}}.
\]
一般地，
\[
Z_{\mathrm{cl}}
=
\frac1h
\sqrt{\frac{2\pi m}{\beta}}
\int\dd x\,\ee^{-\beta V(x)}.
\]
这也说明：构型空间中的 Boltzmann 积分可以通过引入 Gaussian 动量转化为 Hamilton 动力学可采样的相空间积分。

\chapter{Wigner 相空间表象}

\section{Weyl 符号与 Wigner 函数}

对算符 $\hat A$，定义其 Weyl 符号
\[
A_{\Wig}(\vct R,\vct P)
=
\int\dd\vct\Delta\,
\ee^{-\ii\vct P\cdot\vct\Delta/\hbar}
\bra{\vct R+\vct\Delta/2}
\hat A
\ket{\vct R-\vct\Delta/2}.
\]
密度算符 $\hat\rho$ 对应的 Wigner 函数定义为
\[
W_\rho(\vct R,\vct P)
=
\frac{1}{(2\pi\hbar)^F}\rho_{\Wig}(\vct R,\vct P).
\]
对于纯态 $\hat\rho=\ket{\phi}\bra{\phi}$，
\[
W_\phi(\vct R,\vct P)
=
\frac{1}{(2\pi\hbar)^F}
\int\dd\vct\Delta\,
\ee^{-\ii\vct P\cdot\vct\Delta/\hbar}
\phi\!\left(\vct R+\frac{\vct\Delta}{2}\right)
\phi^*\!\left(\vct R-\frac{\vct\Delta}{2}\right).
\]

其边缘分布为
\[
\int\dd\vct P\,W_\phi(\vct R,\vct P)
=
|\phi(\vct R)|^2,
\]
\[
\int\dd\vct R\,W_\phi(\vct R,\vct P)
=
|\widetilde\phi(\vct P)|^2.
\]
Wigner 函数是准概率分布：归一化且具有正确边缘分布，但可取负值。

\section{迹与相空间积分}

Weyl 变换满足
\[
\boxed{
\Tr(\hat A\hat B)
=
\int\frac{\dd\vct R\,\dd\vct P}{(2\pi\hbar)^F}
A_{\Wig}(\vct R,\vct P)
B_{\Wig}(\vct R,\vct P)
}.
\]
因此
\[
\avg{\hat A}
=
\int\dd\vct R\,\dd\vct P\,
W_\rho(\vct R,\vct P)
A_{\Wig}(\vct R,\vct P).
\]

例如，位置算符的 Weyl 符号为
\[
(\hat x)_{\Wig}(R,P)=R.
\]
热密度算符是混态：
\[
\ee^{-\beta\hat H}
=
\sum_n
\ee^{-\beta E_n}
\ket{\phi_n}\bra{\phi_n}.
\]

\section{谐振子的热 Wigner 分布}

令
\[
H_{\mathrm{cl}}(x,p)
=
\frac{p^2}{2m}
+\frac12m\omega^2x^2,
\qquad
a=\frac{\beta\hbar\omega}{2}.
\]
未归一化的热算符 Weyl 符号为
\[
\left[\ee^{-\beta\hat H}\right]_{\Wig}
=
\operatorname{sech}(a)
\exp\left[
-\frac{2\tanh(a)}{\hbar\omega}
H_{\mathrm{cl}}(x,p)
\right].
\]
由此得到
\[
Z
=
\frac{1}{2\sinh a},
\]
以及归一化的热 Wigner 函数
\[
W_\beta(x,p)
=
\frac{\tanh a}{\pi\hbar}
\exp\left[
-\frac{2\tanh a}{\hbar\omega}
H_{\mathrm{cl}}(x,p)
\right].
\]
位置方差为
\[
\avg{x^2}_{\mathrm q}
=
\frac{\hbar}{2m\omega}\coth a.
\]
经典结果为
\[
\avg{x^2}_{\mathrm{cl}}
=
\frac{1}{\beta m\omega^2}.
\]
两者之比
\[
\boxed{
Q
=
\frac{\avg{x^2}_{\mathrm q}}
{\avg{x^2}_{\mathrm{cl}}}
=
\frac{\beta\hbar\omega/2}
{\tanh(\beta\hbar\omega/2)}
}.
\]
高温或经典极限 $\beta\hbar\omega\ll1$ 时，
\[
Q\to1;
\]
低温时量子零点涨落使 $Q>1$。

\section{无界与受约束相空间}

核坐标与动量通常定义在无界的 canonical 相空间上。有限维电子态则可以使用离散 Wigner 表象、spin mapping 或基于 $SU(N)$ coherent state 的受约束相空间。非绝热动力学中，核相空间与电子映射变量相结合，可在统一的相空间框架中描述核--电子耦合。

\chapter{逻辑总结与后续方向}

\begin{routebox}
\[
\begin{aligned}
&\text{经典 Hamilton 动力学}
&&\Rightarrow\text{相空间、能量守恒、Liouville 定理},\\
&\text{统计力学}
&&\Rightarrow\text{系综平均、遍历性、重要性采样},\\
&\text{Lagrange 作用量}
&&\Rightarrow\text{经典驻相路径},\\
&\text{量子传播子}
&&\xrightarrow{\,t=-\ii\hbar\beta\,}\text{Boltzmann 算符},\\
&\text{Trotter 拆分}
&&\Rightarrow\text{虚时间路径积分},\\
&\text{周期虚时间路径}
&&\Rightarrow\text{环形高分子同构与 PIMD},\\
&\text{Weyl 变换}
&&\Rightarrow\text{量子相空间与 Wigner 准概率}.
\end{aligned}
\]
\end{routebox}

后续可沿三条方向继续展开：

\begin{enumerate}
\item \textbf{非绝热动力学：}BO 近似失效时，电子态与核自由度必须耦合处理，包括锥形交叉、无辐射跃迁及光化学过程。
\item \textbf{路径积分估计量与采样：}研究能量、热容、动量分布及相关函数的低方差估计量，以及 normal-mode、staging、Langevin 等高效采样方法。
\item \textbf{量子相空间动力学：}由 Wigner--Moyal 方程出发，理解经典 Liouville 方程的极限，并连接 semiclassical、mapping-variable 与量子轨迹方法。
\end{enumerate}

\end{document}